\section{The temporary geometric Deligne--Illusie boundary}

\begin{definition}[The actual finite exterior-degree source]
\label{def:kn-actual-di-source}
\lean{EnvelopeProperty.Vanishing.DeligneIllusieInput.twistStructure,EnvelopeProperty.Vanishing.DeligneIllusieInput.summandNat,EnvelopeProperty.Vanishing.DeligneIllusieInput.sourceNat,EnvelopeProperty.Vanishing.DeligneIllusieInput.source,EnvelopeProperty.Vanishing.DeligneIllusieInput.realizationFunctor}
\leanok
Let $f:X\to\operatorname{Spec}k$ be a scheme over a field of prime
characteristic $p$, let $X'$ be its actual relative Frobenius twist,
and give $X'$ the structural map defined by the second projection.
For $n\geq0$ take the finite direct sum of the original sheaves
$\Omega^i_{X'/k}$, each placed in its own natural cochain degree $i$,
for $0\leq i\leq n$.  Extend this complex by zero to integer degrees.
For an invertible $L'$ on $X'$, its realized source is obtained by
termwise tensoring with $L'$, forgetting to abelian sheaves, and
localizing in their derived category.  This construction does not
require a decomposition theorem.
\end{definition}

\begin{lemma}[Exact additive tensoring by an actual invertible line]
\label{lem:kn-line-tensor-exact}
\lean{EnvelopeProperty.Vanishing.CartierConnection.tensorFunctor,EnvelopeProperty.Vanishing.CartierConnection.tensorFunctor_isEquivalence,EnvelopeProperty.Vanishing.CartierConnection.tensorFunctor_additive}
\leanok
For an actual invertible module $L'$ on a scheme, the original
functor $M\mapsto M\otimes L'$ is an additive equivalence of the
actual module category, hence is exact.
\end{lemma}
\begin{proof}
Tensoring by the actual dual of $L'$ is an inverse: evaluation,
tensor symmetry and associativity give the two original inverse
comparisons with the tensor unit.  Every equivalence preserves zero
objects, binary products and coproducts.  The biproduct expression
for the sum of two morphisms therefore shows that this original
functor is additive.  Preservation of finite limits and colimits
then gives exactness.
\end{proof}

\begin{lemma}[Extraction of the actual exterior-degree summands]
\label{lem:kn-di-exterior-summand-extraction}
\lean{EnvelopeProperty.Vanishing.DeligneIllusieInput.summandInclusionNat,EnvelopeProperty.Vanishing.DeligneIllusieInput.summandProjectionNat,EnvelopeProperty.Vanishing.DeligneIllusieInput.summandInclusion_projection_nat,EnvelopeProperty.Vanishing.DeligneIllusieInput.realizedSummandIso,EnvelopeProperty.Vanishing.DeligneIllusieInput.realized_summand_retract,EnvelopeProperty.Vanishing.DeligneIllusieInput.subsingleton_realized_summand_hom,EnvelopeProperty.Vanishing.DeligneIllusieInput.subsingleton_hodge_summand_hom,EnvelopeProperty.Vanishing.DeligneIllusieInput.subsingleton_hodge_H}
\leanok
\uses{def:kn-actual-di-source,lem:kn-line-tensor-exact}
For every $0\leq i\leq n$, the original single sheaf
$(\Omega_{X'/k}^i\otimes L')[-i]$ is a retract of the realized
line-twisted decomposition source.  Thus vanishing of the latter's
derived Hom from any fixed object forces vanishing of the derived
Hom into this actual exterior-degree summand.  Taking the source to
be the constant integer sheaf in total degree $t$, the actual
displaced-single/Ext comparison gives the corresponding native
$H^{t-i}(X',\Omega_{X'/k}^i\otimes L')=0$ whenever $t-i\geq0$.
\end{lemma}
\begin{proof}
In the finite coproduct defining the source, take its original
inclusion and the projection that is the identity on the chosen
summand and zero on every other summand.  Their composite on the
chosen summand is the identity.  Extension by zero, tensoring,
forgetting scalars and localization are functors, so they preserve
this exact retraction equation.  Extension of a single complex is
the single complex in the embedded degree; termwise functors also
commute with single complexes.  These canonical comparisons identify
the realized chosen summand with precisely
$(\Omega_{X'/k}^i\otimes L')[-i]$.  If two morphisms into that summand
have equal composites with its inclusion, compose with its
projection to obtain their equality.  When the Hom into the whole
source is zero, their inclusion composites are equal, proving the
asserted vanishing.  For the constant integer sheaf in total degree
$t$, shift both single complexes by $t$.  The source becomes the
single constant sheaf in degree zero and the target becomes the
single original exterior sheaf shifted by $t-i$.  The defining
single-complex/Ext comparison therefore identifies this Hom group
with $\operatorname{Ext}^{t-i}(\mathbb Z,
\Omega_{X'/k}^i\otimes L')$.  This is the native sheaf cohomology
group by definition, proving its vanishing as well.
\end{proof}



\begin{definition}[The actual full decomposition input]
\label{def:kn-geometric-di-input}
\lean{EnvelopeProperty.Vanishing.DeligneIllusieInput.target,EnvelopeProperty.Vanishing.DeligneIllusieInput.Decomposition,EnvelopeProperty.Vanishing.DeligneIllusieInput.WittDecompositionInput,EnvelopeProperty.Vanishing.DeligneIllusieInput.ofWittLift}
\leanok
\uses{def:kn-actual-di-source,lem:kn-witt-special-fibre-dimension}
Let $k$ be a perfect field of prime characteristic $p$.  For a smooth
scheme $f:X\to\operatorname{Spec}k$ of relative dimension $n$, put
$X'=X\times_{\operatorname{Spec}k,F_k}\operatorname{Spec}k$, and let
$F:X\to X'$ be the actual relative Frobenius.  The decomposition is an
isomorphism in $D(\mathcal O_{X'})$
\[
 \bigoplus_{i=0}^{n}\Omega_{X'/k}^{i}[-i]
 \;\simeq\;F_*\Omega_{X/k}^{\bullet}.
\]
The temporary input asserts this isomorphism for the actual special
fibre of every separated smooth relative-dimension-$n$ scheme over $W_2(k)$,
when $n<p$.  It is an explicit hypothesis, not an axiom or an instance.
Given this input and a genuine Witt model, its application chooses the
asserted isomorphism on precisely that model's actual special fibre.
\end{definition}

\begin{lemma}[The genuine Witt special fibre retains its dimension]
\label{lem:kn-witt-special-fibre-dimension}
\lean{EnvelopeProperty.Vanishing.DeligneIllusieInput.specialFibre,EnvelopeProperty.Vanishing.DeligneIllusieInput.specialFibreStructure,EnvelopeProperty.Vanishing.DeligneIllusieInput.specialFibre_smoothOfRelativeDimension}
\leanok
Let $\widetilde f:\widetilde X\to\operatorname{Spec}W_2(k)$ be
smooth of relative dimension $n$.  The actual pullback
$X=\operatorname{Spec}k\times_{\operatorname{Spec}W_2(k)}\widetilde X$
along Witt reduction, with its first projection to $\operatorname{Spec}k$,
is smooth of relative dimension $n$.
\end{lemma}
\begin{proof}
This square is the categorical pullback of $\widetilde f$ along the
actual Witt reduction morphism.  Smoothness of relative dimension $n$
is stable under base change.  Apply this stability to the flipped
pullback square; its structural map is precisely the indicated first
projection.  Neither a choice of coordinates nor a Cartier comparison
is required.
\end{proof}

\begin{lemma}[Tensoring the actual derived decomposition by a line]
\label{lem:kn-di-line-tensor}
\lean{EnvelopeProperty.Vanishing.DeligneIllusieInput.tensorDecomposition}
\leanok
\uses{def:kn-geometric-di-input,lem:kn-line-tensor-exact}
Let $L'$ be an invertible module on the same $X'$.  A supplied actual
decomposition induces the corresponding isomorphism after tensoring
both complexes with $L'$.
\end{lemma}
\begin{proof}
Tensoring by $L'$ is an exact equivalence, whose inverse is tensoring
by its dual.  Hence it preserves quasi-isomorphisms and induces a
functor of the original derived category.  Apply this functor to the
supplied isomorphism, and use its canonical comparison with termwise
tensoring before localization.  These comparisons identify both
endpoints with the indicated actual tensor complexes.
\end{proof}

\begin{lemma}[Underlying abelian derived decomposition]
\label{lem:kn-di-underlying-abelian}
\lean{EnvelopeProperty.Vanishing.DeligneIllusieInput.sheafSource,EnvelopeProperty.Vanishing.DeligneIllusieInput.sheafTarget,EnvelopeProperty.Vanishing.DeligneIllusieInput.sheafDecomposition,EnvelopeProperty.Vanishing.DeligneIllusieInput.sheafTensorDecomposition}
\leanok
\uses{lem:kn-di-line-tensor}
The original module-derived decomposition, also after line tensoring,
induces the corresponding decomposition of the original underlying
abelian-sheaf complexes.
\end{lemma}
\begin{proof}
The underlying-abelian-sheaf functor is additive and preserves finite
limits and colimits, hence is exact.  It induces a functor on the
derived categories.  Apply it to the actual module-derived
isomorphism and compose with the localization comparison at both
endpoints.  These comparisons retain the original underlying sheaf
of every term and the original differential.
\end{proof}

\begin{lemma}[Transport of every original complex differential]
\label{lem:kn-complex-term-transport}
\lean{EnvelopeProperty.Vanishing.CartierConnectionTransport.complex,EnvelopeProperty.Vanishing.CartierConnectionTransport.iso,EnvelopeProperty.Vanishing.CartierConnectionTransport.differential_comp}
\leanok
Let $K$ be a complex in a category with zero morphisms and let
$e_i:K^i\simeq A_i$ be actual isomorphisms.  On the terms $A_i$,
the maps $e_jd_{ij}e_i^{-1}$ form a complex of the same shape,
isomorphic to the whole original complex through the original $e_i$.
\end{lemma}
\begin{proof}
Every map outside the original shape is zero because the corresponding
$d_{ij}$ is zero.  A composite of two transported differentials is
$e_kd_{jk}e_j^{-1}e_jd_{ij}e_i^{-1}$; cancellation of the middle
inverse pair reduces it to $e_kd_{jk}d_{ij}e_i^{-1}=0$.
The identity $e_jd_{ij}=(e_jd_{ij}e_i^{-1})e_i$ proves compatibility
with every original differential.  The inverse term isomorphisms give
the inverse complex morphism, so this is a complex isomorphism.
\end{proof}

\begin{definition}[The actual line-twisted Cartier complex]
\label{def:kn-cartier-connection-complex}
\lean{EnvelopeProperty.Vanishing.CartierConnection.twistedComplex,EnvelopeProperty.Vanishing.CartierConnection.termProjectionIso,EnvelopeProperty.Vanishing.CartierConnection.connectionTerm,EnvelopeProperty.Vanishing.CartierConnection.connectionDifferential,EnvelopeProperty.Vanishing.CartierConnection.connectionComplex}
\leanok
Let $L'$ be invertible on $X'$.  Tensor the $\mathcal O_{X'}$-linear
complex $F_*\Omega_X^\bullet$ by $L'$.  The line projection formula
identifies its degree-$i$ term with
$F_*(\Omega_X^i\otimes F^*L')$.  Transporting the differentials
through these actual term isomorphisms defines the corresponding
Cartier-connection complex on these pushed modules.
\end{definition}

\begin{lemma}[The transported connection retains the original complex]
\label{lem:kn-cartier-connection-complex-iso}
\lean{EnvelopeProperty.Vanishing.CartierConnection.connectionDifferential_square,EnvelopeProperty.Vanishing.CartierConnection.connectionComplexIso}
\leanok
\uses{def:kn-cartier-connection-complex,lem:kn-complex-term-transport}
The term projection formulas form an isomorphism from the original
line-twisted relative-Frobenius complex to the Cartier-connection
complex.  In particular its differential has square zero.
\end{lemma}
\begin{proof}
Write $e_i$ for the projection formula in degree $i$ and $d_i$ for
the original twisted differential.  Define the transported
differential as $e_{i+1}d_i e_i^{-1}$.  The composite of two successive
differentials is $e_{i+2}d_{i+1}d_i e_i^{-1}=0$.  The same formula
states exactly that $e_i$ commutes with the differentials; its inverse
also does, giving the required complex isomorphism.
\end{proof}

The ordinary de Rham differential on $X$ is not
$\mathcal O_X$-linear.  The construction tensors only after the
actual relative Frobenius pushforward, on $X'$, where its differential
is linear.  It does not replace the missing projection, field-twist
or hypercohomology comparisons by vanishing assumptions.

\begin{lemma}[The amplified original Hodge terms]
\label{lem:kn-amplified-cartier-terms}
\lean{EnvelopeProperty.Vanishing.CartierConnectionPower.twistLinePower,EnvelopeProperty.Vanishing.CartierConnectionPower.twistLinePower_isInvertible,EnvelopeProperty.Vanishing.CartierConnectionPower.twistedComplex,EnvelopeProperty.Vanishing.CartierConnectionPower.term,EnvelopeProperty.Vanishing.CartierConnectionPower.termIso,EnvelopeProperty.Vanishing.CartierConnectionPower.differential,EnvelopeProperty.Vanishing.CartierConnectionPower.differential_square,EnvelopeProperty.Vanishing.CartierConnectionPower.complex,EnvelopeProperty.Vanishing.CartierConnectionPower.complexIso,EnvelopeProperty.Vanishing.CartierConnectionPower.sheafTermIso}
\leanok
\uses{cor:kn-frobenius-line-projection,lem:kn-cartier-connection-complex-iso}
For an original line $L$ on $X$, let $L'$ be its pullback to the
actual twist.  In every degree $q$ the actual complex
$F_*\Omega_X^\bullet\otimes (L')^m$ has its original term canonically
isomorphic to $F_*(L^{pm}\otimes\Omega_X^q)$.  Transporting the
original differential through these isomorphisms gives an isomorphic
complex on precisely these pushed positive-power modules.
\end{lemma}
\begin{proof}
Pullback preserves invertibility, and the $m$th power of the pulled
line has inverse the $m$th power of its dual, including the tensor
unit when $m=0$.  Thus every line used in this projection formula
is actually invertible.
The actual Frobenius line comparison gives
$F^*((L')^m)\simeq L^{pm}$.  Apply the actual line projection
formula to $\Omega_X^q$ and $(L')^m$, use this pullback comparison,
and use tensor symmetry to put the original power on the left.
This is the indicated term isomorphism.  Transport of the original
differentials preserves square zero and gives a whole complex
isomorphism, by the preceding transport argument.  Extension to
integer degrees has the canonical original term comparison at every
nonnegative degree, and the underlying abelian-sheaf functor
preserves these term isomorphisms.
\end{proof}

\begin{lemma}[The actual amplified target has bounded terms]
\label{lem:kn-amplified-cartier-bounded}
\lean{EnvelopeProperty.Vanishing.CartierConnectionPower.twistedTerm_isZero,EnvelopeProperty.Vanishing.CartierConnectionPower.term_isZero,EnvelopeProperty.Vanishing.CartierConnectionPower.complex_extend_isStrictlyLE,EnvelopeProperty.Vanishing.CartierConnectionPower.sheafComplex,EnvelopeProperty.Vanishing.CartierConnectionPower.sheafComplex_X_isZero_below,EnvelopeProperty.Vanishing.CartierConnectionPower.sheafComplex_X_isZero_above}
\leanok
\uses{lem:kn-amplified-cartier-terms}
The underlying integer-indexed actual amplified complex is zero in
every degree below zero and in every degree above $n$.
The corresponding natural-degree module terms and the extended
transported complex obey the same upper bound.
\end{lemma}
\begin{proof}
Negative terms are zero by extension from the natural-number-indexed
original de Rham complex.  Above $n$, the original exterior sheaves
are zero: on every genuine smooth affine chart the exterior basis
is indexed by $q$-element subsets of an $n$-element set, hence is
empty for $q>n$.  These chart equalities detect the zero original
sheaf.  Pushforward, tensoring and forgetting scalars all preserve
zero objects.  Thus the same zero bounds hold for the actual
amplified underlying complex.
\end{proof}

\begin{lemma}[The decomposition has the actual amplified target]
\label{lem:kn-di-amplified-target}
\lean{EnvelopeProperty.Vanishing.CartierConnectionPower.realizedSource,EnvelopeProperty.Vanishing.CartierConnectionPower.decompositionIso,EnvelopeProperty.Vanishing.CartierConnectionPower.source_hodge_H}
\leanok
\uses{lem:kn-di-underlying-abelian,lem:kn-amplified-cartier-terms,lem:kn-di-exterior-summand-extraction}
The supplied original Deligne--Illusie isomorphism gives a derived
isomorphism from its realized source tensored by $(L')^m$ to the
actual amplified underlying-abelian complex used above.
If the realized source has zero derived Hom from the constant integer
sheaf in total degree $t$, then for each $0\leq i\leq n$ with
$t-i\geq0$ its native group
$H^{t-i}(X',\Omega^i_{X'/k}\otimes(L')^m)$ is zero.
\end{lemma}
\begin{proof}
The original line remains invertible after pullback and taking the
$m$th tensor power.  Apply the exact line-tensor and underlying
abelian-sheaf derived functors to the supplied original
decomposition.  The target is, by its definition, precisely the
underlying abelian complex obtained by tensoring the extended
original de Rham complex with $(L')^m$.  No change of original
differential, scheme or coefficient projection is made.
Apply the original summand retractions and the native single/Ext
comparison to obtain the last assertion.  The displayed derived
isomorphism also transports zero Hom from the actual target to the
same realized source, so it can be combined directly with bounded
hypercohomology vanishing.
\end{proof}
